\ExplSyntaxOn

\int_new:N \l__binp_int
\int_set:Nn \l__binp_int {4}
\tl_new:N \l__bina_tl
\tl_new:N \l__binb_tl
\tl_new:N \l__binhw_tl
\int_new:N \l__bina_int
\int_new:N \l__binb_int
\int_new:N \l__binc_int
\int_new:N \l__bind_int
\fp_new:N \l__bina_fp

\DeclareDocumentCommand \BinaryPrecision {m}
{
  \int_set:Nn \l__binp_int {#1}
}

\cs_new_nopar:Npn \int_to_bin:Nn #1#2
{
  \tl_clear:N #1
  \int_set:Nn \l__bina_int {#2}
  \prg_replicate:nn {\l__binp_int}
  {
    \int_set:Nn \l__binb_int {\int_mod:nn {\l__bina_int} {2}}
    \tl_put_left:Nx #1 {\int_use:N \l__binb_int}
    \int_set:Nn \l__bina_int {\int_div_truncate:nn {\l__bina_int} {2}}
  }
}

\cs_new_nopar:Npn \bin_hamming_weight:NN #1#2
{
  \tl_set_eq:NN \l__binhw_tl #2
  \tl_replace_all:Nnn \l__binhw_tl {0} {}
  \int_set:Nn #1 {\tl_count:N \l__binhw_tl}
}

\cs_new_nopar:Npn \bin_flip_one:NNn #1#2#3
{
  \tl_set_eq:NN #2#1
  \prg_replicate:nn {#3 - 1}
  {
    \tl_replace_once:Nnn #2 {1} {a}
  }
  \tl_replace_once:Nnn #2 {1} {0}
  \tl_replace_all:Nnn #2 {a} {1}
}

\DeclareDocumentCommand \HammingWeight {m m}
{
  \int_to_bin:Nn \l__binhw_tl {#2}
  \bin_hamming_weight:NN #1 \l__binhw_tl
}

\DeclareDocumentCommand \HyperCube {m}
{
  \int_set:Nn \l__binp_int {#1}
  \int_zero_new:c {l__hyper_0_int}
  \int_set:cn {l__hyper_0_int} {-1}
  \int_step_inline:nnnn {1} {1} {#1}
  {
    \int_zero_new:c {l__hyper_##1_int}
    \int_set:cn  {l__hyper_##1_int} {\int_use:c {l__hyper_\int_eval:n {##1 - 1}_int} * (#1 - ##1 + 1) / ##1}
  }
  \fp_set:Nn \l__bina_fp {2^{#1}-1}
  \int_set:Nn \l__binc_int {\fp_to_int:N \l__bina_fp}
  \int_step_inline:nnnn {0} {1} {\l__binc_int}
  {
    \int_to_bin:Nn \l__bina_tl {##1}
    \bin_hamming_weight:NN \l__binb_int \l__bina_tl
    \node[every~ hypercube~ label/.try] (hg- \l__bina_tl) at (\int_use:c {l__hyper_\int_use:N \l__binb_int _int},\int_use:N \l__binb_int) {\texttt{\l__bina_tl}};
    \int_incr:c {l__hyper_\int_use:N \l__binb_int _int}
    \int_incr:c {l__hyper_\int_use:N \l__binb_int _int}
  }
  \int_step_inline:nnnn {0} {1} {\l__binc_int}
    {
      \int_to_bin:Nn \l__bina_tl {##1}
      \bin_hamming_weight:NN \l__binb_int \l__bina_tl
      \int_step_inline:nnnn {1} {1} {\l__binp_int}
      {
        \bin_flip_one:NNn \l__bina_tl \l__binb_tl {####1}
        \tl_if_eq:NNTF \l__bina_tl \l__binb_tl
          { % do nothing if nodes are the same
          }
          {
            \draw[every~ hypercube~ edge/.try] (hg- \l__bina_tl) -- (hg- \l__binb_tl);
          }
      }
    }
  }

\ExplSyntaxOff

% Macro to cross out a single node
\newcommand{\CrossOutNode}[1]{%
  \draw[red, thick]
    ($(#1) + (-0.3, 0.3)$)
    --
    ($(#1) + (0.3, -0.3)$);
  \draw[red, thick]
    ($(#1) + (-0.3, -0.3)$)
    --
    ($(#1) + (0.3, 0.3)$);
}

% Macro to cross out a list of nodes
\newcommand{\CrossOutNodes}[1]{%
  \foreach \n in {#1} {%
    \CrossOutNode{\n}%
  }
}
% Macro to draw a circle around a single node
\newcommand{\CircleNode}[2][RoyalBlue]{%
  \draw[#1, thick]
    ($(#2)$)
    circle (0.5);
}

% Macro to circle a list of nodes
\newcommand{\CircleNodes}[2][RoyalBlue]{%
  \foreach \n in {#2} {%
    \CircleNode[#1]{\n}%
  }
}

% Überwachungstabelle der Aufgabe mit 10 Nachrichtenbits und 4 Kontrollbits
% (Spalten 1 bis 14, Zeilen c1 bis c4); Massstab im tikzpicture setzen
\newcommand{\ueberwachungZehnVier}{%
  \fill[blue!25] (0,0) rectangle (15,1);
  \fill[blue!15] (0,-4) rectangle (1,0);
  \fill[gray!8] (1,-4) rectangle (15,0);
  \foreach \i in {1,...,10} {
      \node at (\i+.5,.5) {\(a_{\i}\)};
    }
  \foreach \i in {1,...,4} {
      \pgfmathtruncatemacro{\col}{10+\i}
      \node[text=orange!80!black] at (\col+.5,.5) {\(c_{\i}\)};
      \node[text=orange!80!black] at (.5,.5-\i) {\(c_{\i}\)};
    }
  % Überwachungsbereiche von c1 bis c4, eine Liste pro Zeile
  \foreach \cols [count=\z] in {{1,2,3,5,6,7,11},{1,2,4,5,8,9,12},{1,3,4,6,8,10,13},{2,3,4,7,9,10,14}} {
      \foreach \col in \cols {
          \node at (\col+.5,.5-\z) {\(\times\)};
        }
    }
  % Trennung Nachrichtenbits / Kontrollbits
  \draw[blue!70!black,very thick] (11,-4) -- (11,1);
  \draw[step=1,black!30,thin] (0,-4) grid (15,1);
  \draw[black!65] (0,-4) rectangle (15,1);
}

% Bitwerte des Codeworts über den Spaltentiteln
% #1: 14 Bits als Komma-Liste, #2: Spalte magenta hervorheben (0 = keine)
\newcommand{\ueberwachungWerte}[2]{%
  \node[font=\small\itshape,anchor=east] at (1,1.5) {Wert};
  \foreach \w [count=\col] in {#1} {
      \ifnum\col=#2
        \node[magenta,font=\small\bfseries\ttfamily] at (\col+.5,1.5) {\w};
      \else
        \node[font=\small\ttfamily] at (\col+.5,1.5) {\w};
      \fi
    }
}

% Anzahl Einsen unter den Kreuzen jeder Zeile rechts der Tabelle,
% ungerade Anzahlen (Kontrollbit meldet) magenta
\newcommand{\ueberwachungEinsen}[1]{%
  \node[font=\small\itshape] at (16.4,.5) {Einsen};
  \foreach \n [count=\z] in {#1} {
      \ifodd\n
        \node[magenta] at (16.4,.5-\z) {\textbf{\n}\ \(\boldsymbol{\times}\)};
      \else
        \node[green!50!black] at (16.4,.5-\z) {\n\ \checkmark};
      \fi
    }
}


% command to make sure colorboxes don't create whitespace between code lines
\fboxsep=.6pt

% default settings for listings
\lstset{style=python}

\title[Datenintegrität]{Datenintegrität}

\subtitle{Kartentrick-Codierung}

\author[Blum, Cyril] % (optional)
{Cyril Blum}

\institute[KLW] % (optional)
{
    Fachschaft Informatik\\
    Kantonsschule im Lee
}

\date[\today] % (optional)

\logo{\includegraphics[width=2cm]{Logos/FreeFlower.pdf}}

\titlegraphic{\vspace{-.7cm}\includegraphics[width=.85\textwidth]{"Grundlagen_Info/06_Datenbanken/Figures/title-emojis"}}

\frame{\titlepage}


\begin{frame}{Rückblick 1-fehlerkorrigierende Codierung}
    \begin{itemize}[<+->]
        \item Drei Farben, Codewörter der Länge 5 mit Abstand 3: von Hand lösbar
        \item Nachrichten mit 16 Bit $\rightarrow$ \(2^{16}=65536\) Codewörter
        \item Von Hand suchen ist hoffnungslos
        \item Gesucht: ein Verfahren, das Codewörter systematisch erzeugt
    \end{itemize}
\end{frame}

% Kartenraster für den Zaubertrick: alle Schritte nutzen dasselbe 5x5-Raster,
% damit die 4x4-Karten beim Folienwechsel an ihrem Platz bleiben und nur die
% neue Reihe (oben) und Spalte (rechts) dazukommen.
\newcommand{\zauberx}{1.05}  % horizontaler Abstand der Kartenmitten (cm)
\newcommand{\zaubery}{1.35}  % vertikaler Abstand der Kartenmitten (cm)
% \zauberreihe{Zeile}{Karte § Karte § …}; Zeile 0 ist oben, "null" lässt einen Platz leer
\newcommand{\zauberreihe}[2]{%
    \setsepchar{§}\readlist*\zauberkarten{#2}%
    \foreach \c in {1,...,\zauberkartenlen}{%
        \itemtomacro\zauberkarten[\c]\zauberkarte
        \node[inner sep=0pt] at (\c*\zauberx,-#1*\zaubery) {\AffCarteJeu[Hauteur=1.05,StyleJeu=v1]{\zauberkarte}};
    }%
}
% farbige Hinterlegung der neuen Reihe und Spalte (Kontrollkarten)
\newcommand{\zauberneu}{%
    \foreach \c in {1,...,5} {\fill[orange!35,rounded corners=2pt] (\c*\zauberx,0) +(-.48,-.65) rectangle +(.48,.65);}
    \foreach \r in {1,...,4} {\fill[orange!35,rounded corners=2pt] (5*\zauberx,-\r*\zaubery) +(-.48,-.65) rectangle +(.48,.65);}
}

\begin{frame}{\enquote{Zaubertrick}}
    \begin{center}
        \begin{tikzpicture}
            % festes Raster für alle Schritte
            \useasboundingbox (\zauberx-.48,.65) rectangle (5*\zauberx+.48,-4*\zaubery-.65);
            \only<1>{
                \zauberreihe{1}{DK § Dos § 10T § Dos}
                \zauberreihe{2}{Dos § AP § AT § VP}
                \zauberreihe{3}{VT § Dos § 5C § 3T}
                \zauberreihe{4}{Dos § Dos § 5C § VP}
            }
            \only<2>{
                \zauberneu
                \zauberreihe{0}{Dos § 8C § Dos § VP § Dos}
                \zauberreihe{1}{DK § Dos § 10T § Dos § Dos}
                \zauberreihe{2}{Dos § AP § AT § VP § 8T}
                \zauberreihe{3}{VT § Dos § 5C § 3T § 4P}
                \zauberreihe{4}{Dos § Dos § 5C § VP § Dos}
            }
            \only<3>{
                \zauberneu
                \zauberreihe{0}{Dos § 8C § Dos § VP § Dos}
                \zauberreihe{1}{DK § Dos § 10T § Dos § Dos}
                \zauberreihe{2}{Dos § AP § AT § VP § 8T}
                \zauberreihe{3}{VT § Dos § Dos § 3T § 4P}
                \zauberreihe{4}{Dos § Dos § 5C § VP § Dos}
            }
        \end{tikzpicture}

        \vspace{.5em}
        \only<1>{\textbf{Erster Schritt}: Auslegen der ursprünglichen Karten ($n=m=4$)}%
        \only<2>{\textbf{Zweiter Schritt}: Hinzufügen einer Reihe und einer Spalte (\colorbox{orange!35}{Hilfszauberer})}%
        \only<3>{\textbf{Dritter Schritt}: Karten, nachdem jemand eine Karte umgedreht hat}
    \end{center}
\end{frame}

\begin{frame}{\enquote{Zaubertrick} (binär als Matrix)}
    \only<1>{
        \[
            \begin{pmatrix}
                1 & 0 & 1 & 0 \\
                0 & 1 & 1 & 1 \\
                1 & 0 & 1 & 1 \\
                0 & 0 & 1 & 1 \\
            \end{pmatrix}
        \]
        \vspace{0.5em}
        \textbf{Erster Schritt}: Ursprüngliche Bitmatrix ($n=m=4$)
    }
    \only<2>{
        \[
            \begin{pmatrix}
                \underline{0} & \underline{1} & \underline{0} & \underline{1} & \underline{0} \\
                1 & 0 & 1 & 0 & \underline{0} \\
                0 & 1 & 1 & 1 & \underline{1} \\
                1 & 0 & 1 & 1 & \underline{1} \\
                0 & 0 & 1 & 1 & \underline{0} \\
            \end{pmatrix}
        \]
        \vspace{0.5em}
        \textbf{Zweiter Schritt}: Hinzufügen einer Reihe und einer Spalte (\underline{Kontrollbits})
    }
    \only<3>{
        \[
            \begin{pmatrix}
                \underline{0} & \underline{1} & \underline{0} & \underline{1} & \underline{0} \\
                1 & 0 & 1 & 0 & \underline{0} \\
                0 & 1 & 1 & 1 & \underline{1} \\
                1 & 0 & 0 & 1 & \underline{1} \\
                0 & 0 & 1 & 1 & \underline{0} \\
            \end{pmatrix}
        \]
        \vspace{0.5em}
        \textbf{Dritter Schritt}: Eine Bitkarte wurde umgedreht (Fehlerfall)
    }
\end{frame}


\begin{frame}{Die Kartentrickmethode}
    Für die folgende 6-stellige Bitfolge soll eine 1-fehlerkorrigierende Codierung gefunden werden: \texttt{010101}
    \begin{columns}
        \begin{column}{.45\textwidth}
            \begin{table}[H]
                \centering
                \begin{tabular}{ccc}
                    0 & 1 & 0 \\\midrule
                    1 & 0 & 1 \\
                \end{tabular}
                \caption*{Ursprüngliche Nachricht}
                \label{tab:ex-zauberer}
            \end{table}
        \end{column}
        \begin{column}{.05\textwidth}
            \only<2>{$\rightarrow$}
        \end{column}
        \begin{column}{.45\textwidth}
            \only<2>{
                \begin{table}[H]
                \centering
                \begin{tabular}{cccc}
                    \underline{1} & \underline{1} & \underline{1} & \underline{1}    \\\midrule
                    0             & 1             & 0             & \underline{1} \\\midrule
                    1             & 0             & 1             & \underline{0} \\
                \end{tabular}
                \caption*{Ursprüngliche Nachricht + \underline{Kontrollbits}}
                \label{tab:ex-zauberer-parity}
            \end{table}
            }
        \end{column}
    \end{columns}
    % Reihenfolge im Codewort wie im Skript: erste Zeile, letzte Spalte, zuletzt das Bit oben rechts
    \only<2>{
        \vspace{.5em}
        Codewort: Nachricht, erste Zeile und letzte Spalte je ohne das Bit oben rechts, zuletzt das Bit oben rechts\\[.3em]
        \texttt{0 1 0 1 0 1 \underline{1 1 1} \underline{1 0} \underline{1}}
    }
\end{frame}

\begin{frame}{Auftrag}{Skript \enquote{Datenintegrität}}
    \begin{itemize}
        \item \aufgabebox{1.17, 1.18}
        \item \challengebox{1.16}
    \end{itemize}
\end{frame}

\begin{frame}{Optimalität: so wenig Kontrollbits wie möglich}
    Mehrere Möglichkeiten für Nachrichten der Länge 16
    \begin{itemize}[<+->]
        \item $16 = 4 \cdot 4 \rightarrow$ \underline{Kontrollbits: $4 + 4 + 1 = 9$}
        \item $16 = 2 \cdot 8 \rightarrow$ \underline{Kontrollbits: $2 + 8 + 1 = 11$}
        \item $16 = 1 \cdot 16 \rightarrow$ \underline{Kontrollbits: $1 + 16 + 1 = 18$}
        \item Die beiden Seiten des Rechtecks sollten möglichst gleich gross sein
        \item Bei Nachrichten der Länge 48 $\rightarrow$ $6 \cdot 8$ ist optimal!
    \end{itemize}
\end{frame}
